\section{Synth Modules}\label{synth-modules}

A closer look at the synthesis modules used across the voices.

\subsection{Oscillator}\label{oscillator}

The main oscillator of a voice has an \textbf{oscillator type}
(\texttt{typ}) that decides how the \texttt{wav} selection is read.

{\def\LTcaptype{none} % do not increment counter
\begin{longtable}[]{@{}
  >{\raggedright\arraybackslash}p{(\linewidth - 4\tabcolsep) * \real{0.3333}}
  >{\raggedright\arraybackslash}p{(\linewidth - 4\tabcolsep) * \real{0.3333}}
  >{\raggedright\arraybackslash}p{(\linewidth - 4\tabcolsep) * \real{0.3333}}@{}}
\toprule\noalign{}
\begin{minipage}[b]{\linewidth}\raggedright
\texttt{typ}
\end{minipage} & \begin{minipage}[b]{\linewidth}\raggedright
Name
\end{minipage} & \begin{minipage}[b]{\linewidth}\raggedright
Plays
\end{minipage} \\
\midrule\noalign{}
\endhead
\bottomrule\noalign{}
\endlastfoot
\texttt{wav} & Stock waves & \texttt{Sin}, \texttt{Tri}, \texttt{Saw},
\texttt{Rec} (rectangle), \texttt{Noi} (noise). \\
\texttt{smp} & One-shot sample & A sample from the installed catalogue,
played once from the trigger. \\
\texttt{slp} & Looped sample & A sample played in its loop region. \\
\texttt{wtb} & Wavetable & A wavetable with frame blending; \texttt{pos}
selects the frame. \\
\texttt{pwm} & Pulse-width modulation & Sin, Tri, Saw or Rec with a
variable pulse width set by \texttt{pos}. \\
\end{longtable}
}

Under \texttt{smp}, \texttt{slp} and \texttt{wtb}, \texttt{wav} values
from the sixth entry upward select installed catalogue entries, named
\texttt{s01} onward (samples), \texttt{l01} onward (looped samples) and
\texttt{w01} onward (wavetables). Which kind of entry plays is decided
by \texttt{typ}.

\verify{The built-in crash sample no longer exists in flash. A file \texttt{02\_crash.wav} on the card replaces it. Confirm the shipped factory sample set.}

\subsubsection{\texorpdfstring{Position
(\texttt{pos})}{Position (pos)}}\label{position-pos}

\texttt{pos} is the fourth cell of the first oscillator screen on every
voice. Its meaning follows \texttt{typ}:

{\def\LTcaptype{none} % do not increment counter
\begin{longtable}[]{@{}
  >{\raggedright\arraybackslash}p{(\linewidth - 2\tabcolsep) * \real{0.5000}}
  >{\raggedright\arraybackslash}p{(\linewidth - 2\tabcolsep) * \real{0.5000}}@{}}
\toprule\noalign{}
\begin{minipage}[b]{\linewidth}\raggedright
Type
\end{minipage} & \begin{minipage}[b]{\linewidth}\raggedright
Meaning of \texttt{pos} (0 to 127)
\end{minipage} \\
\midrule\noalign{}
\endhead
\bottomrule\noalign{}
\endlastfoot
\texttt{wav} & Start phase: \texttt{pos}/128 of a cycle at each trigger.
Noise ignores it. A move while the note plays glides within one
block. \\
\texttt{smp} & Start the sample at \texttt{pos}/128 of its length. \\
\texttt{slp} & Start at \texttt{pos}/128 of the sample, wrapped into the
loop. \\
\texttt{wtb} & Frame position in the wavetable. \\
\texttt{pwm} & Pulse width: 0 is 50 percent (the unmodified wave), 127
is the narrowest pulse. \\
\end{longtable}
}

\subsubsection{\texorpdfstring{Run mode
(\texttt{run})}{Run mode (run)}}\label{run-mode-run}

\begin{paramtable}
rst & Reset      & Each retrigger resets the carrier phase. This is the traditional behaviour. \\
fre & Free       & The carrier phase is never reset. Modulators slew, the amplitude envelope is gain-matched and the filter keeps ringing. \\
itp & Interpolate & The phase is reset and the old waveform is crossfaded out over 132 samples. A fully decayed \texttt{itp} equals \texttt{rst}. \\
\end{paramtable}

\subsubsection{Wavetable interpolation}\label{wavetable-interpolation}

When a wave-modulated oscillator blends between neighbouring waveforms
or wavetable frames, the firmware interpolates for at most two
oscillators at a time; additional ones step. Slowly modulated
oscillators are interpolated; fast ones may step.

\subsection{Samples and Wavetables}\label{samples-and-wavetables}

Samples and wavetables are installed from the SD card into the flash
memory with \textbf{Load \textgreater{} Samples} (see the Storage
chapter).

\begin{itemize}
\tightlist
\item
  Accepted files are \textbf{mono, 16-bit, 44.1 kHz PCM WAV}.
\item
  Put samples in \texttt{/samples}, and wavetables in
  \texttt{/wavetables}, one sub-folder per wavetable.
\item
  Names are sorted in natural order (\texttt{2} before \texttt{10}). On
  the display a long name shows its first four and last four characters,
  with any numeric prefix removed.
\item
  Every wavetable frame is resampled to 1,024 points at install time.
  8-bit files are not installed.
\item
  Samples are kept through power-down. Installing erases and rewrites
  the sample area.
\end{itemize}

\attention{Installing samples erases the sample area first. A card prepared with an older firmware
shows the catalogue as stale and must be reinstalled before the samples can be used.}

\figplaceholder{sample_folders.png}{Folder layout for samples and wavetables.}{Tree view of the SD card root showing samples, wavetables with two example wavetable folders, and which names appear on the display as s01, l01 and w01.}

\subsection{Filter}\label{filter}

The voice filter is a two-pole (12 dB) state-variable filter realised as
a non-linear zero-delay-feedback model. It shapes the harmonic content:
a hi-hat needs a highpass to remove low frequencies, a clap benefits
from a bandpass.

Three parameters control it: frequency, resonance and drive.

\begin{itemize}
\tightlist
\item
  \textbf{Frequency} sets the operating point.
\item
  \textbf{Resonance} controls the feedback path and boosts frequencies
  around the cutoff.
\item
  \textbf{Drive} controls how hard the filter is driven. At 0 it is
  clean and nearly linear; higher values engage a soft clipper and the
  slew limit of the integrators, so peaks distort first and then the
  whole signal.
\end{itemize}

\subsubsection{Filter types}\label{filter-types}

\figplaceholder{filter_lowpass.png}{Lowpass filter response.}{Magnitude response with three resonance settings: low in red, medium in green, high in blue. Same cutoff in every plot, marked by a vertical dashed line. Use image\_templates/filter\_plot\_template.svg.}

\figplaceholder{filter_highpass.png}{Highpass filter response.}{As the lowpass plot, for the highpass type.}

\figplaceholder{filter_bandpass.png}{Bandpass filter response.}{As the lowpass plot, for the bandpass type.}

\figplaceholder{filter_unitgain_bandpass.png}{Unity-gain bandpass response.}{As the lowpass plot; gain is always 0 dB at the peak and resonance controls only the width of the passband.}

\figplaceholder{filter_notch.png}{Notch filter response.}{As the lowpass plot, for the notch type; resonance controls the width of the stopband.}

\figplaceholder{filter_peak.png}{Peak filter response.}{As the lowpass plot, for the peak type; resonance controls the amount of boost.}

\begin{itemize}
\tightlist
\item
  \textbf{Lowpass} removes high frequencies.
\item
  \textbf{Highpass} removes low frequencies.
\item
  \textbf{Bandpass} removes both sides of the cutoff.
\item
  \textbf{Unity-gain bandpass} is a bandpass scaled to 0 dB; resonance
  controls the width.
\item
  \textbf{Notch} removes frequencies around the cutoff.
\item
  \textbf{Peak} boosts frequencies around the cutoff.
\item
  \textbf{Naive two-pole low-pass} is a cheaper low-pass kept for kick
  transients.
\end{itemize}

The Alkaline voice uses a different filter, a \textbf{diode ladder} with
its own resonance and drive laws; see the Instruments chapter.

\subsection{Envelopes}\label{envelopes}

Whenever a voice is triggered, its envelopes restart. The signal rises
at the speed set by attack, reaches its maximum and falls at the speed
set by decay.

\figplaceholder{envelope_linear.png}{Linear envelope curve.}{Amplitude against time showing attack and decay with slope 63 (a straight-line decay). Use image\_templates/envelope\_template.svg.}

\begin{itemize}
\tightlist
\item
  \textbf{Attack and decay} control the duration of each stage. Higher
  values mean longer times. With attack at zero the envelope starts
  directly in the decay stage.
\item
  \textbf{Slope} controls the curve. 63 is linear, lower values are
  exponential and higher values logarithmic.
\end{itemize}

\figplaceholder{envelope_exponential.png}{Low slope: exponential decay.}{As the linear envelope with a low slope value.}

\figplaceholder{envelope_logarithmic.png}{High slope: logarithmic decay.}{As the linear envelope with a high slope value.}

\begin{itemize}
\tightlist
\item
  \textbf{Repeat} simulates clap-like attacks. With a repeat count above
  zero, the attack stage is replaced by a repeat stage in which the
  envelope retriggers rapidly; attack sets the duration of the whole
  repeat phase.
\end{itemize}

\figplaceholder{envelope_repeat.png}{Envelope in repeat mode.}{Envelope with repeat count 3, repeat phase in blue and decay phase in red.}

Each voice also has a \textbf{pitch envelope} (a one-shot decay) that
sweeps the oscillator pitch. On the Drum voice it also scales the depth
of the FM, sync and ring modes.

\subsection{Transient Generator}\label{transient-generator}

The transient generator shapes the first milliseconds of a sound.

\begin{itemize}
\tightlist
\item
  \textbf{Snappy (\texttt{Snp}):} an additional fast pitch envelope on
  the oscillator produces a snap or click. \texttt{vol} is the depth and
  \texttt{frq} the decay.
\item
  \textbf{Offset (\texttt{Ofs}):} adjusts the start phase of the
  oscillators. At volume 0 they start at a zero crossing; at maximum
  volume the wave starts at its peak, giving a loud pop.
\item
  \textbf{Click samples:} twelve built-in clicks, and then every
  installed sample (\texttt{s01}\ldots) as a one-shot. \texttt{frq} sets
  the playback speed (1.0 is the file's own rate); \texttt{vol} is
  applied on top.
\end{itemize}

\subsection{Wavefolder}\label{wavefolder}

With the voice chain set to \textbf{WFd} the transient stage is replaced
by a \textbf{wavefolder} that sits before the filter. The \texttt{wav}
cell selects the fold type, \texttt{vol} its depth and \texttt{frq} a
secondary control that depends on the type.

\begin{paramtable}
Tri & Triangle fold    & Classic triangle fold. \\
TrD & Triangle, decaying & Fold depth decays over time (5 ms to 2 s). \\
TrI & Triangle, intense & Deeper fold range. \\
Chb & Chebyshev        & Polynomial fold with unity small-signal gain. \\
Fd2 & Two-stage fold   & Two cascaded folds; \texttt{frq} sets the second stage ratio. \\
Sin & Sine fold        & Smooth sine folding. \\
SnD & Sine, decaying   & Sine fold with decaying depth. \\
Asn & Sine with accent & Sine fold with an even-harmonic accent. \\
\end{paramtable}

\verify{Fold type names and order on screen, and the exact meaning of frq for each type.}

At depth 0 the triangle folds do nothing.

\subsection{Distortion and Sample-Rate
Reduction}\label{distortion-and-sample-rate-reduction}

\begin{itemize}
\tightlist
\item
  \textbf{Distortion} is a soft saturator that is always in the signal
  path, with drive 0 acting as a nearly clean setting.
\item
  \textbf{Sample-rate reduction} (\texttt{srt} or \texttt{sr}) is
  sample-and-hold decimation. It is applied only to the voices that have
  it in their chain (Drum, Snare, Cymbal and HiHat). Alkaline does not
  have it.
\item
  There is \textbf{no global sample-rate setting}; the former global
  control was replaced by the Effect Morph (\texttt{fxm}) cell on the
  PERF page.
\end{itemize}

\subsection{LFO}\label{lfo}

See the LFO page in the Instruments chapter. Further notes:

\begin{itemize}
\tightlist
\item
  \textbf{Rate table:} rate 0 stops the LFO, 1 is a 100 second period,
  and the upper half of the range moves in semitone steps so that LFO
  rates relate musically.
\item
  \textbf{Synced rates} run from 4/1 down to 1/32 and follow the
  sequencer tempo (internal or external).
\item
  An LFO can modulate any modulatable parameter of any voice, the
  Scene's Morph amounts (\texttt{scn}) and the Effect parameters
  (\texttt{fx}). It modulates \emph{around} the value set by the menu,
  Morph or automation, and the base is restored when the target is
  cleared.
\end{itemize}

\subsection{Velocity Modulation}\label{velocity-modulation}

Each voice routes note velocity to one destination parameter with an
amount. Velocity cannot modulate Effect parameters. On Alkaline,
velocity is the accent.
